\subsection{Recommender system}\label{subsec:recommendation-methodology}

We rank the canonical microbial active-substance value. It represents a BCA or a listed BCA mixture. We keep
products as evidence so users can inspect the authorisation records behind each candidate.

We build the current ranking catalogue from Portuguese authorisation records only. Portugal is the only
source that provides complete records for the required fields. We keep the other country outputs separate until we
recover and validate their missing links.

We derive the interaction table from accepted \textit{final} rows; \textit{review} rows remain a separate
unresolved queue. We normalise microbial mixture strings into a canonical order and remove duplicate components.
Each row represents one crop-pathogen context and one canonical BCA value. We collapse repeated source rows into
an interaction count and aggregate source, output class, product name, and product function. Support
uses distinct products when available and interaction counts otherwise. Products remain \textit{evidence}, not ranked items.

We treat each crop-pathogen context as an implicit user and each BCA as an item. The interaction table becomes a
context-item matrix. We retain interaction counts for support features but use binary presence for similarity.
The collaborative component computes cosine similarity between BCA vectors, prefers contexts with the same
pathogen, and falls back to global contexts when necessary. It shrinks similarities when item overlap is small
and scores each candidate from its two strongest similarities to the observed query BCAs. If the query has no
observed items, the collaborative score is zero and direct support provides the fallback.

The ranker combines context, pathogen, and collaborative support. It also computes crop, total, and dataset
support, but assigns them zero weight. It normalises count features within the candidate set before weighting them.
The current weights are $0.475$ for known context
matches, $0.190$ for context support, $0.285$ for pathogen support, and $0.050$ for collaborative similarity;
crop, total, and dataset support have weight zero. The query also reports the triage score as a diagnostic, but
does not use it for ordering in the current evaluation. Its weight is zero and the value remains available for
audit and manual review.

The query matches case-insensitive literal text in the scientific crop and common or scientific pathogen fields.
It separates known matches from inferred candidates. Inferred candidates come from related authorisation evidence;
they do not establish a new authorisation or a biological efficacy claim.